% This is samplepaper.tex, a sample chapter demonstrating the
% LLNCS macro package for Springer Computer Science proceedings;
% Version 2.21 of 2022/01/12
%
\documentclass{llncs}
%
\usepackage[T1]{fontenc}
% T1 fonts will be used to generate the final print and online PDFs,
% so please use T1 fonts in your manuscript whenever possible.
% Other font encondings may result in incorrect characters.
%
\usepackage{graphicx}
% Used for displaying a sample figure. If possible, figure files should
% be included in EPS format.
%
% If you use the hyperref package, please uncomment the following two lines
% to display URLs in blue roman font according to Springer's eBook style:
%\usepackage{color}
%\renewcommand\UrlFont{\color{blue}\rmfamily}
%\urlstyle{rm}
%

\usepackage[style=apa]{biblatex} % <- AGREGAR ESTA LÍNEA
\addbibresource{thebiblio.bib} % <- AGREGAR ESTA LÍNEA
\begin{document}
%
% ENGLISH VERSION
\title{Anomaly Detection in biomedical data by means of Autoencoders}

%\titlerunning{Abbreviated paper title}
% If the paper title is too long for the running head, you can set
% an abbreviated paper title here
%
%\author{Matías Cisnero\inst{1} \and
%Juliana Gambini\inst{1,2}\orcidID{0000-0002-4534-1402}}
%%
%\authorrunning{M Cisnero et al.}
%% First names are abbreviated in the running head.
%% If there are more than two authors, 'et al.' is used.
%%
%\institute{
%LIDEC, Universidad Nacional de Hurlingham, Pcia. de Buenos Aires, Argentina\\
%\email{juliana.gambini@unahur.edu.ar} \\
%\and
%CPSI, Universidad Tecnológica Nacional, Regional Buenos Aires}
%%
\maketitle              % typeset the header of the contribution
%
\begin{abstract}
This line of research addresses the detection of anomalies in biomedical data using the Autoencoder method.

Autoencoders are artificial neural networks with a symmetrical architecture in which the input layer is the same as output layer. 
This method reduces the dimension of original data and then reconstructs it. The reconstruction error is the key point of this research.

The hypothesis of this work is that patients who present a large reconstruction error (greater than a threshold) can be classified as anomalous, which can help the medical professional to detect diseases.

With this objective, the Autoencoder method is implemented
and trained with different proportions of data from healthy and sick patients. Several architectures are used to optimize the parameters.

Three variants are evaluated: the classic autoencoder, the sparse autoencoder, and the contractive autoencoder.
The threshold used to classify these data as normal or anomalies based on their reconstruction error is computed using the Youden index.

In this case, the methods are tested using the
Breast Cancer Wisconsin (Diagnostic) dataset, which has samples from patients with benign and malignant patterns obtained by fine-needle aspiration (FNA).

The results are encouraging, but further research is needed to obtain more conclusive results.

% keywords in lowercase
\keywords{Anomaly detection\and Biomedical Data\and Autoencoders}
\end{abstract}

% SPANISH VERSION
% Clear the previous title info and set Spanish version

% Spanish title and abstract
\begin{center}
\vspace{1cm}
{\Large \bfseries\boldmath Detección de Anomalías en datos biomédicos utilizando Autoencoders}

\vspace{0.5cm}
\end{center}

\begin{abstract}
La presente línea de investigación aborda la detección de anomalías en datos biomédicos utilizando el método Autoencoder.

 Los Autoencoders son redes neuronales artificiales que poseen una arquitectura simétrica, en la que la capa de entrada es exactamente igual a la de salida. Este método reduce la dimensión de los datos originales para luego reconstruirlos. El error de reconstrucción es la pieza clave de esta investigación. 

La hipótesis de este trabajo es que los pacientes  que presenten un error de reconstrucción grande (mayor que un umbral), pueden ser clasificados como anómalos, lo cual puede ayudar al profesional médico a detectar enfermedades.  

Con este objetivo, se implementa el método Autoencoder
y se entrena con distintas proporciones de datos de pacientes sanos y enfermos. Se utilizan diferentes arquitecturas para optimizar los parámetros. 



Se evalúan tres variantes: el autoencoder clásico, el
\emph{sparse} autoencoder y el \emph{contractive} autoencoder.
 El umbral utilizado para clasificar estos datos como normales o anomalías según su
error de reconstrucción, se calcula utilizando el índice de Youden.

En este caso, los métodos se prueban utilizando el conjunto de datos
Breast Cancer Wisconsin (Diagnostic), que posee muestras de pacientes con patrones benignos y malignos obtenidas mediante punción con aguja fina (FNA). 

Los resultados son alentadores pero es necesario profundizar la investigación para obtener resultados más contundentes. 
% palabras clave en minúsculas
\keywords{ Detección de Anomalías\and  Datos biomédicos\and Autoencoders}
\end{abstract}


\vspace{0.5cm}
\setcounter{page}{1}

%
%\section{First Section}
%\subsection{A Subsection Sample}
%Please note that the first paragraph of a section or subsection is
%not indented. The first paragraph that follows a table, figure,
%equation etc. does not need an indent, either.
%
%Subsequent paragraphs, however, are indented.
%
%\subsubsection{Sample Heading (Third Level)} Only two levels of
%headings should be numbered. Lower level headings remain unnumbered;
%they are formatted as run-in headings.
%
%\paragraph{Sample Heading (Fourth Level)}
%The contribution should contain no more than four levels of
%headings. Table~\ref{tab1} gives a summary of all heading levels.
%
%\begin{table}
%\caption{Table captions should be placed above the
%tables.}\label{tab1}
%\begin{tabular}{|l|l|l|}
%\hline
%Heading level &  Example & Font size and style\\
%\hline
%Title (centered) &  {\Large\bfseries Lecture Notes} & 14 point, bold\\
%1st-level heading &  {\large\bfseries 1 Introduction} & 12 point, bold\\
%2nd-level heading & {\bfseries 2.1 Printing Area} & 10 point, bold\\
%3rd-level heading & {\bfseries Run-in Heading in Bold.} Text follows & 10 point, bold\\
%4th-level heading & {\itshape Lowest Level Heading.} Text follows & 10 point, italic\\
%\hline
%\end{tabular}
%\end{table}
%
%
%\noindent Displayed equations are centered and set on a separate
%line.
%\begin{equation}
%x + y = z
%\end{equation}
%Please try to avoid rasterized images for line-art diagrams and
%schemas. Whenever possible, use vector graphics instead (see
%Fig.~\ref{fig1}).
%
%\begin{figure}
%\includegraphics[width=\textwidth]{fig1.eps}
%\caption{A figure caption is always placed below the illustration.
%Please note that short captions are centered, while long ones are
%justified by the macro package automatically.} \label{fig1}
%\end{figure}
%
%\begin{theorem}
%This is a sample theorem. The run-in heading is set in bold, while
%the following text appears in italics. Definitions, lemmas,
%propositions, and corollaries are styled the same way.
%\end{theorem}
%%
%% the environments 'definition', 'lemma', 'proposition', 'corollary',
%% 'remark', and 'example' are defined in the LLNCS documentclass as well.
%%
%\begin{proof}
%Proofs, examples, and remarks have the initial word in italics,
%while the following text appears in normal font.
%\end{proof}
%For citations of references, we use the APA 7 convention.  The following bibliography provides
%a sample reference list with entries for journal
%articles~\parencite{ref_article1}, an LNCS chapter~\parencite{ref_lncs1}, a
%book~\parencite{ref_book1}, proceedings without editors~\parencite{ref_proc1},
%and a homepage~\parencite{ref_url1}. Multiple citations are grouped
%\parencite{ref_article1,ref_lncs1,ref_book1},
%\parencite{ref_article1,ref_book1,ref_proc1,ref_url1}.
%
%\begin{credits}
%\subsubsection{\ackname} A bold run-in heading in small font size at the end of the paper is
%used for general acknowledgments, for example: This study was funded
%by X (grant number Y).
%
%\subsubsection{\discintname}
%It is now necessary to declare any competing interests or to specifically
%state that the authors have no competing interests. Please place the
%statement with a bold run-in heading in small font size beneath the
%(optional) acknowledgments,
%for example: The authors have no competing interests to declare that are
%relevant to the content of this article. Or: Author A has received research
%grants from Company W. Author B has received a speaker honorarium from
%Company X and owns stock in Company Y. Author C is a member of committee Z.
%\end{credits}
%%
%% ---- Bibliography ----
%%

%\printbibliography

\end{document}